% \documentclass[paper=a4, fontsize=11pt]{scrartcl} % A4 paper and 11pt font size
\documentclass[11pt, a4paper]{book}
\usepackage[T1]{fontenc} % Use 8-bit encoding that has 256 glyphs
\usepackage[utf8]{inputenc}
\usepackage{fourier} % Use the Adobe Utopia font for the document - comment this line to return to the LaTeX default
\usepackage{listings} % para insertar código con formato similar al editor
\usepackage[spanish, es-tabla]{babel} % Selecciona el español para palabras introducidas automáticamente, p.ej. "septiembre" en la fecha y especifica que se use la palabra Tabla en vez de Cuadro
\usepackage{url} % ,href} %para incluir URLs e hipervínculos dentro del texto (aunque hay que instalar href)
\usepackage{graphics,graphicx, float} %para incluir imágenes y colocarlas
\usepackage[gen]{eurosym} %para incluir el símbolo del euro
\usepackage{cite} %para incluir citas del archivo <nombre>.bib
\usepackage{enumerate}
\usepackage{hyperref}
\usepackage{graphicx}
\usepackage{tabularx}
\usepackage{booktabs}

\usepackage[table,xcdraw]{xcolor}
\hypersetup{
	colorlinks=true,	% false: boxed links; true: colored links
	linkcolor=black,	% color of internal links
	urlcolor=cyan		% color of external links
}

\renewcommand{\familydefault}{\sfdefault}
\usepackage{fancyhdr} % Custom headers and footers
\pagestyle{fancyplain} % Makes all pages in the document conform to the custom headers and footers
\fancyhead[L]{} % Empty left header
\fancyhead[C]{} % Empty center header
\fancyhead[R]{Ángel Píñar Rivas} % My name
\fancyfoot[L]{} % Empty left footer
\fancyfoot[C]{} % Empty center footer
\fancyfoot[R]{\thepage} % Page numbering for right footer
%\renewcommand{\headrulewidth}{0pt} % Remove header underlines
\renewcommand{\footrulewidth}{0pt} % Remove footer underlines
\setlength{\headheight}{13.6pt} % Customize the height of the header

\usepackage{titlesec, blindtext, color}
\definecolor{gray75}{gray}{0.75}
\newcommand{\hsp}{\hspace{20pt}}
\titleformat{\chapter}[hang]{\Huge\bfseries}{\thechapter\hsp\textcolor{gray75}{|}\hsp}{0pt}{\Huge\bfseries}
\setcounter{secnumdepth}{4}
\usepackage[Lenny]{fncychap}


\begin{document}
	\begin{titlepage}
\newlength{\centeroffset}
\setlength{\centeroffset}{-0.5\oddsidemargin}
\addtolength{\centeroffset}{0.5\evensidemargin}
\thispagestyle{empty}

\noindent\hspace*{\centeroffset}\begin{minipage}{\textwidth}

\centering
\includegraphics[width=0.9\textwidth]{logos/logo_ugr.jpg}\\[1.4cm]

\textsc{ \Large TRABAJO FIN DE MÁSTER\\[0.2cm]}
\textsc{ GRADO EN INGENIERÍA INFORMÁTICA}\\[1cm]

{\Huge\bfseries My Many Reads \\}
\noindent\rule[-1ex]{\textwidth}{3pt}\\[3.5ex]
{\large\bfseries Plataforma social para lectores con sistema de recomendaciones }
\end{minipage}

\vspace{2.5cm}
\noindent\hspace*{\centeroffset}
\begin{minipage}{\textwidth}
\centering

\textbf{Autor}\\ {Ángel Píñar Rivas}\\[2.5ex]
\textbf{Director}\\ {Juan Julían Merelo Guervós}\\[2cm]
\includegraphics[width=0.3\textwidth]{logos/etsiit_logo.png}\\[0.1cm]
\textsc{Escuela Técnica Superior de Ingenierías Informática y de Telecomunicación}\\
\textsc{---}\\
Granada, Septiembre de 2022
\end{minipage}
\end{titlepage}

	\thispagestyle{empty}

\begin{center}
{\large\bfseries My Many Reads \\ Plataforma social para lectores con sistema de recomendaciones }\\
\end{center}
\begin{center}
Ángel Píñar Rivas\\
\end{center}

%\vspace{0.7cm}

\vspace{0.5cm}
\noindent{\textbf{Palabras clave}: \textit{software libre}, \textit{desarrollo ágil}, \textit{integración continua}, \textit{proceso de desarrollo}, \textit{control de calidad}, \textit{biblioteca}, \textit{gestión}, \textit{lecturas}, \textit{libros}, \textit{recomendaciones}, \textit{estadísticas}, \textit{popularidad}, \textit{tendencias}, \textit{usuarios}
\vspace{0.7cm}

\noindent{\textbf{Resumen}\\

El objetivo de este proyecto es definir un proceso de desarrollo claro y robusto, apoyándose en una mentalidad ágil y un conjunto de herramientas que garanticen la buena calidad del código y su documentación. Para poner estas prácticas en uso, se desarrollará un sistema de gestión de lecturas que ofrezca recomendaciones automáticas basadas en las lecturas de cada usuario además de la posibilidad de participar activamente en la plataforma creando recomendaciones personalizadas y otro contenido que pueda resultar de interés al resto de usuarios. Mediante la información y estadísticas de uso generadas en esta plataforma, se obtendrá el beneficio necesario para cubrir costes de mantenimiento.

\cleardoublepage

\begin{center}
	{\large\bfseries My Many Reads \\ Social platform for readers with recommendation system }\\
\end{center}
\begin{center}
	Ángel Píñar Rivas\\
\end{center}
\vspace{0.5cm}
\noindent{\textbf{Keywords}: \textit{open source}, \textit{agile development}, \textit{continuous integration}, \textit{development process}, \textit{quality assurance}, \textit{library}, \textit{management}, \textit{readings}, \textit{books}, \textit{recommendations}, \textit{statistics}, \textit{popularity}, \textit{trends}, \textit{users}
\vspace{0.7cm}

\noindent{\textbf{Abstract}\\

The main goal of this project is to define a clear and robust development process, relying on an agile mindset and a set of tools to ensure good quality of both code and documentation. To put these practices into use, a readings management system will be developed to offer automatic recommendations based on the readings of each user as well as the possibility to actively participate in the platform by creating personalized recommendations and other content that may be of interest to the rest of users. Through the information and usage statistics generated on this platform, the necessary profit will be obtained to cover maintenance costs.

\cleardoublepage

\thispagestyle{empty}

\noindent\rule[-1ex]{\textwidth}{2pt}\\[4.5ex]

D. \textbf{Juan Julián Merelo Guervós}, Profesor del Departamento de Arquitectura y Tecnología de Computadores de la Universidad de Granada.

\vspace{0.5cm}

\textbf{Informo:}

\vspace{0.5cm}

Que el presente trabajo, titulado \textit{\textbf{My Many Reads}},
ha sido realizado bajo mi supervisión por \textbf{Ángel Píñar Rivas}, y autorizo la defensa de dicho trabajo ante el tribunal
que corresponda.

\vspace{0.5cm}

Y para que conste, expiden y firman el presente informe en Granada a Septiembre de 2022.

\vspace{1cm}

\textbf{El director: }

\vspace{5cm}

\noindent \textbf{(Juan Julián Merelo Guervós)}

\chapter*{Agradecimientos}

A la resistencia, por el apoyo frente a las oleadas incesantes de presión y trabajo durante este máster, y la sensación de compañerismo en tiempos de pandemia. Sin vosotros, la sensación de soledad me habría llevado a abandonar este curso.

A Phoenix, por su empatía, su disposición a escuchar y la sensación de espalda cubierta, todo ello sin obligación de hacerlo.

A mi tutor, por aguantar mis cabezonerías con mi propia forma de trabajo, las revisiones frecuentes y las discusiones fructíferas que me han enseñado tanto información como puntos de vista.

A mis amigos y familia por escuchar y convertir mi agobio y preocupaciones en risas y diversión, cuidando de mi salud mental.

Y finalmente a mí mismo (aunque suene feo), por resistir la tentación de abandonar y seguir adelante hasta el final, aún con todo el agobio y estrés, y las motivaciones desaparecidas.




	\newpage
	\tableofcontents

	\newpage
	\listoffigures

	\listoftables 
	\newpage

	\chapter{Introducción}

A día de hoy, el afán por la lectura está aumentando entre no sólo los jóvenes (el sector que más preocupaba) sino adultos también. Una de las posibles causas de este hecho podría haber sido la pandemia, ya que con la cuarentena muchos han tenido que encontrar nuevos \textit{hobbies} a los que dedicarse desde su hogar para reemplazar actividades al aire libre.

Como con cualquier otro tipo de afición, no es sencillo iniciarse en el mundo de la lectura sin conocer un buen punto de partida, pues lo ideal es comenzar con algo fácilmente disfrutable, lo cual depende fuertemente de los gustos personales de la persona en concreto, por lo que en general se suele hacer una pequeña investigación sobre aquellos que podrían ser interesantes.

Este tipo de búsquedas no solo las realizan aquellas personas introduciéndose en la lectura, sino también los ya acostumbrados lectores que acaban con su lista de libros pendientes o prefieren descubrir libros nuevos. Al fin y al cabo, cualquier persona interesada en la lectura acaba buscando algo que leer por los medios de que disponga.

\section{Motivación}

Justo antes de la pandemia decidí explorar nuevos \textit{hobbies} que me separaran de la pantalla. Compré varios juegos de mesa, empecé a pintar miniaturas y decidí comprarme un libro con una premisa medianamente interesante. De estos tres, las miniaturas fueron las únicas que acabé descartando, porque aunque me gustan, son bastante caras, y como para los juegos de mesa dependes de organizarte con tus amigos, me enfoqué bastante en la lectura.

Me encontré con el problema de que no conocía a mucha gente que leyera, y además, que compartiera mis gustos, por lo que dependía de buscar en internet en base a lo poco que llevaba leído, y la verdad es que era una tarea algo tediosa. Se me ocurrió que podría existir una plataforma de \textit{reviews} de la que podría servirme, y me topé con \textit{Goodreads}, pero la verdad es que el sitio no terminaba de gustarme del todo. Conozco otras plataformas para series y películas que sí son bastante cómodas de usar para estas búsquedas, pero no para libros.

Finalmente me acabé acostumbrando bastante a mi método "manual" de buscar nuevos libros, introduje a un par de amigos con mis mismos gustos a la lectura para discutir con ellos y acabé disfrutando más del \textit{hobby} en general. Sin embargo, siempre pensé (y sigo pensando) que todo esto habría sido más fácil e incluso mejor, si hubiera una plataforma online en la que explorar libros, compartir opiniones y recomendaciones, y gestionar las lecturas, en cierto modo parecido a lo que ofrece \textit{Goodreads} pero eliminando todos los detalles, aspectos y funcionalidades que hacen que yo, a pesar de disfrutar mucho de la lectura y de llevarme bien con los ordenadores en general, no quiera usar esta plataforma que parece estar pensada para gente con mi misma afición, y que además es la más destacable de este ámbito con diferencia.

\section{Descripción del problema}

Existe una gran diversidad de géneros y gustos personales a tener en cuenta a la hora de realizar una buena elección de un libro. En el caso de no tener contactos que puedan aconsejar, habría que recurrir a otras fuentes de información, principalmente internet ya que es la más accesible, amplia y fácil de utilizar de las disponibles.

Se pueden encontrar ideas o sugerencias con simples búsquedas a través de Google, pero éstas suelen venir de revistas o blogs cuyos motivos de dichas recomendaciones no son transparentes, pudiendo ser, por ejemplo, totalmente personales del autor (una única persona) o haberse visto influenciados por un movimiento de patrocinio por parte de alguna editorial.

Esto no parece una solución ideal, pues un usuario promedio fácilmente puede desconfiar de estas recomendaciones, ya que por una parte se desconoce de dónde proceden y cuales son los motivos detrás de éstas, y por otra parte podrían no encajar con los gustos del usuario.

El propósito de este proyecto es acabar con este problema, permitiendo a todo aquel usuario con interés en la lectura servirse de unas recomendaciones totalmente personalizadas, apoyándose por opiniones que una multitud de usuarios similares aportan para cada libro. En otras palabras, elaborando las recomendaciones para los usuarios en base a sus lecturas y respaldando esas recomendaciones con opiniones de otros usuarios reales, por lo que se acaba creando una plataforma social en torno a la lectura en la que se refleja la opinión mayoritaria de sus usuarios, aportando mucha fiabilidad.

\section{Objetivos}
En el caso de \textit{My Many Reads}, el objetivo es construir una plataforma social en la que crear una comunidad de usuarios alrededor de la lectura. Partiendo de una base de datos de libros, se pretende reunir a una comunidad de usuarios con intereses en la lectura, ofreciéndoles una serie de servicios y funcionalidades como pueden ser:

\begin{itemize}
    \item Gestión de una biblioteca personal, en la que los usuarios puedan almacenar los libros leídos, pendientes y abandonados junto a una valoración.
    \item Sistema de recomendaciones automático construyendo los gustos de los usuarios a partir de sus bibliotecas personales.
    \item Sistema de búsqueda de libros por popularidad en la plataforma o sus características propias como pueden ser géneros, idioma o fecha de publicación.
    \item Foro de discusiones para cada libro en el que plasmar opiniones, teorías u otros pensamientos relacionados con dicho libro.
    \item Posibilidad de escribir recomendaciones personalizadas entre dos libros similares, que puedan ayudar a otros usuarios a tomar una decisión con su próxima lectura.
\end{itemize}

La meta global por tanto es la de construir una plataforma con características atractivas para atraer a un conjunto de usuarios. Sabiendo esto, el siguiente paso es establecer una lista de objetivos con el fin de planificar y diseñar en torno al cumplimiento de estos, enfocando todo el desarrollo en una dirección. En el caso de \textit{My Many Reads}, estos son los objetivos:

\begin{itemize}
    \item Cumplir con una funcionalidad mínima de consulta de libros y gestión listas o bibliotecas.
    \item Con el fin de que los usuarios puedran encontrar fácilmente buenas sugerencias para sus próximas lecturas, construir un sistema de recomendaciones que, a partir de los gustos del usuario, devuelva una lista de libros que podrían ser de interés para éste.
    \item Permitir a los usuarios participar activamente en creación de contenido de la plataforma mediante foros de discusiones, recomendaciones entre libros personalizadas y reviews.
    \item Servir como medio de comunicación para ofrecer noticias pertenecientes al mundo de la lectura.
    \item Recoger estadísticas de los usuarios de la plataforma para conocer información sobre, por ejemplo, las tendencias actuales.
\end{itemize}

Estos dos últimos objetivos tienen cierta dependencia a que el proyecto se encuentre desplegado y con una base mínima de usuarios. Para llegar a ese punto, es imprescindible cumplir los objetivos previos cuya intención es atraer usuarios a la plataforma.

	\chapter{Descripción del problema}

	\chapter{Estado del arte}

Actualmente solo destacan otras dos plataformas similares, Goodreads\cite{Goodreads} y Librarything\cite{Librarything}, siendo la primera la opción más conocida y popular en cuanto a este tipo de plataformas. A continuación se exploran las características de estas dos plataformas, listando aspectos positivos y negativos.

\section{Goodreads}

Sin duda la opción más popular, esta plataforma ofrece una serie de opciones bastante llamativas y útiles:

\begin{itemize}
    \item Una biblioteca de libros a la que puedes añadir libros leídos, en progreso y en lista de espera, junto con tu puntuación, fechas de inicio y fin de lectura y \textit{review}.
    \item Un sistema de recomendaciones basado en tu biblioteca.
    \item Estadísticas de tus lecturas.
    \item Listas creadas por la comunidad que recogen libros con características comunes, como puede ser \textit{Books That Should Be Made Into Movies}.
    \item Un foro de discusiones, un apartado de citas famosas y secciones de AMA (\textit{Ask Me Anything}) dirigidas a autores.
    \item Una sección de noticias y artículos relacionados con los libros y el mundo de la lectura en general.
    \item Elementos de \textit{gamificación} de la lectura como preguntas diarias, un \textit{challenge} de lecturas anuales y cuestionarios acerca de libros.
    \item Permite a los usuarios escribir y compartir sus propias historias.
\end{itemize}

No por nada esta plataforma es la más popular, pero sin embargo no carece de aspectos mejorables:

\begin{itemize}
    \item La interfaz en sí es bastante mejorable, se hace difícil a un usuario promedio navegar y descubrir todo lo que ofrece sin tener que realizar un poco de esfuerzo.
    \item No permite explorar cómodamente la base de datos de libros, solo puedes ver las tendencias actuales (decididas por la plataforma mediante algún algoritmo desconocido) y explorar de un género concreto.
    \item Las listas mencionadas anteriormente las crea un único usuario, por lo que solo la opinión de una persona es la que dicta la relación entre todos los libros de la lista.
    \item El apartado de discusión de libros no ofrece mucha libertad, una vez lo encuentras parece más bien un conjunto de publicaciones filtrado según si son acerca de o relacionados con el libro concreto.
    \item Al visitar la página de un libro, si quieres consultar opiniones de la gente, lo que se te ofrece es su puntación media (no precisamente muy visible) y una sección de comentarios como la que te podrías encontrar en Amazon o cualquier tienda online. El foro de discusiones mencionado anteriormente es una fuente valiosa de \textit{feedback} que, aún con los problemas mencionados anteriormente, no parece ser accesible desde esta página.
    \item Información adicional como la existencia de adaptaciones en otros formatos (tales como películas o series) no aparece por ninguna parte.
    \item No se hace mención a diferentes ediciones de un libro ni a los idiomas en los que está disponible.
    \item El autor, otro elemento del cual un usuario podría usar para guiarse a la hora de buscar su siguiente lectura, solo dispone de un perfil muy similar al que podría tener cualquier usuario común, diferenciándose en una sección de AMA y de listado de sus libros. No parece haber un lugar en el que discutir sobre el estilo narrativo u otros elementos comunes a todos sus libros, que puedan llevar a un usuario a leer diferentes libros del autor que no compartan género entre sí.
\end{itemize}

\section{Librarything}

La opción menos popular, no da indicios de la funcionalidad al completo que ofrece hasta que el usuario se registra. Además de ofrecer menos funcionalidad que Goodreads\cite{Goodreads}, parece que tiene algunos problemas adicionales. Comenzando por lo que ofrece esta plataforma:

\begin{itemize}
\item Una biblioteca de libros para cada usuario, junto con un sistema de recomendaciones automáticas en base a esta biblioteca.
\item Una sección de comunidades, similares a grupos de lectura, en las que compartir opiniones, embarcarse en retos o hacer concursos.
\item Un foro de discusiones, similar al de Goodreads, pero más accesible.
\item Búsqueda de autores.
\item Conexión de tu perfil con otras aplicaciones, como redes sociales o plataformas de intercambio de libros.
\item Muestra más metadatos sobre los libros que Goodreads\cite{Goodreads}, como idiomas en los que está disponibles y portadas que puedes encontrar.
\item Permite consultar estadísticas de los usuarios de la plataforma, lo cual resulta de utilidad a la hora de decidir si involucrarse con esta plataforma.
\end{itemize}

La funcionalidad ofrecida es claramente menor, aunque cabe decir que tiene lo suficiente como para poder atraer a un conjunto de usuarios bastante amplio. Pero de nuevo, no todo se ofrece de la mejor forma posible, existen aspectos susceptibles de mejora:

\begin{itemize}
\item De nuevo, la interfaz resulta claramente mejorable, lo que puede causar pérdida de usuarios.
\item Se usan diferentes bibliotecas según el estado de lectura de un libro, y solo permite al usuario añadir una puntuación y etiquetas a cada libro añadido a la biblioteca.
\item No existe un buscador libros o base de datos propia, lo que permite agregar dos veces el mismo libro si se busca en páginas o ediciones diferentes, por lo que aparecen inconsistencias.
\item Tiene un buscador de librerías o bibliotecas cercanas a tu zona, lo cual no está mal, pero existen servicios similares más populares y fiables, como el propio Google Maps.
\item Como en Goodreads\cite{Goodreads}, solo permite explorar libros por género o etiqueta, sin ningún tipo de filtro más allá de eso.
\item No existen listas de libros y las recomendaciones escritas por usuarios no tienen por qué incluir el motivo de la recomendación, lo cual no genera mucha confianza.
\end{itemize}

\section{Conclusiones}

Aunque también hay que mencionar que existen otras plataformas de este estilo pero con temática diferente, en cuanto a libros hay pocas opciones, y la que parece la mejor de ellas, Goodreads\cite{Goodreads}, tiene una amplitud de aspectos que se podrían mejorar, y por tanto introducirse en este entorno no parece una lucha imposible.
	\chapter{Planificación}

\section{Metodología utilizada}

\section{Temporización}

\section{Seguimiento del desarrollo}

	\chapter{Análisis del problema}

\section{Aprovechamiento de la plataforma}
Si bien es cierto que ya existe algunas plataformas similares (ver [estado del arte]()), se puede ver que no son perfectas, y de alguna manera se podrían introducir mejoras en algún sentido que podrían aumentar el número de usuarios, por lo que no sería descabellado pensar en hacerles competencia y tener algo de éxito.
\\\\
El primer paso sería lograr una comunidad activa, para lo cual hay que ofrecer incentivos a los usuarios para emplear la plataforma, pues una red social sin usuarios no serviría para nada. Pero también hay que tener en cuenta que ofrecer un mero servicio a los usuarios sin obtener algún tipo de beneficio con el que cubrir los gastos de mantenimiento no lo más adecuado, por lo que habría que explotar este tráfico de usuarios e información de alguna manera:
\\\\
\begin{itemize}
    \item Sirviendo como lugar publicitario para libros de todo tipo, ya que en la plataforma se concentrarían usuarios con interés en la lectura, y por tanto, posibles clientes.
    \item Así mismo, se podrían recoger datos sobre los gustos e intereses de los usuarios para que las propias editoriales sepan qué géneros y autores están más de moda, ayudándoles así a embarcarse en proyectos con mayor beneficio potencial.
\end{itemize} 

	\chapter{Implementación}

	\chapter{Conclusiones, estado actual y futuro del proyecto}

En este último capítulo se recogerán las conclusiones obtenidas tras el desarrollo del proyecto hasta este punto, se resumirá el estado actual, con todas las características implementadas y se comentarán los pasos a seguir en etapas futuras de desarrollo.

\section{Conclusiones}

No existe una razón concreta por la que yo quise realizar un proyecto relacionado con libros, sino más bien mi intención fue la de poner a prueba mi capacidad para organizar un proyecto, así como emplear tecnologías que en mi trabajo no uso, con el fin de aprender un poco más y salir de mi zona de comfort.

Gracias a esto he aprendido bastante, y no solo con lo finalmente implementado en el proyecto sino con las numerosas alternativas que se consideraron y se pusieron a prueba en pequeños PoC (\textit{Proof of Concept}). Esto me ha permitido hacer una pequeña visita al estado de algunas tecnologías que se usan actualmente y me ha permitido contrastarlo con mi conocimiento previo, por lo que ahora me veo con una mejor capacidad de toma de decisiones al haber explorado más a fondo nuevas opciones.

También me ha permitido reafirmar algunas sospechas sobre ciertas formas de organizar el trabajo. Comencé pensando, erróneamente, que al ser un único desarrollador, el tamaño de las epopeyas y tareas refinadas debía más pequeño. Tras un primer intento de aplicar esta filosofía me dí cuenta de que los costes en tiempo de la gestión administrativa de estas tareas empezaban a ser bastante altos en proporción al tiempo necesario para resolver dichas tareas, es decir, programar.

Por otra parte, los beneficios de esta gestión administrativa tienen menos impacto al tratarse de un trabajo en solitario. Al ser un único desarrollador, lo más común es tener el conocimiento y los objetivos de cada historia bastante claros, y parte de la intención de refinar y explicar apropiadamente cada una de las tareas a desarrollar es que otros desarrolladores puedan fácilmente trabajar en ellas sin tener que buscar información adicional del contexto de la tarea.

También me ha sorprendido la utilidad que tiene, a nivel de revisión, abrir \textit{pull requests} cuando solo hay un desarrollador. Aunque es cierto que mi tutor los revisaba, también me daba tiempo a ``aparcar'' la tarea un tiempo y desconectar, permitiéndome así revisar mi propio trabajo.

En cuanto a tecnologías, he aprendido bastante al estudiar y poner en práctica muchas de estas a las cuales no estoy acostumbrado. Un ejemplo de esto son los \textit{ORM}, de los cuales no conocía su existencia, y me parece una opción muy buena para determinados tipos de aplicaciones. Por otra parte, \textit{python} ha sido desde siempre un lenguaje que nunca me ha terminado de gustar como ya expliqué al momento de tomar esta decisión, pero al ver la cantidad de complementación que un \textit{linter} puede aportarle, como por ejemplo las anotaciones de tipo, ha hecho que mi opinión mejore, aunque sigue teniendo características que no me terminan de gustar como la complejidad de la gestión de los entornos de desarrollo.

En resumen, este proyecto me ha permitido aprender sobre diferentes tipos de bases de datos, la implementación de sistemas de integración continua, de los cuales solo me había servido de su utilidad anteriormente, y de la importancia de la configuración del entorno de desarrollo para acomodar el trabajo. Por otra parte, he podido experimentar el cierre del servicio de la base de datos que estaba usando, lo cual me ha permitido experimentar de primera mano que la migración es un proceso realmente simple, y que aunque estas situaciones resulten molestas, lidiar con ellas acaba siendo bastante sencillo.

\section{Estado actual}

\section{El futuro del proyecto}
	
	\newpage
	\bibliography{biblio}
	\bibliographystyle{plain}
	
\end{document}

